\subsection{Resampling, Bootstrap, and Permutation Methods}

Resampling methods approximate sampling distributions by recomputing statistics on data-generated pseudo-samples.

\begin{definition} \label{def:nonparametric-bootstrap} \hlt{Nonparametric Bootstrap}\\
Given data $X_1,\ldots,X_n$, the \hlt{nonparametric bootstrap} samples $X_1^\ast,\ldots,X_n^\ast$ independently from the empirical measure $\widehat{P}_n$.  The conditional distribution of
\[
T(\widehat{P}_n^\ast)-T(\widehat{P}_n)
\]
given the data is used to approximate the sampling distribution of $T(\widehat{P}_n)-T(P)$.
\end{definition}

\begin{method} \label{meth:bootstrap-standard-error} \hlt{Bootstrap Standard Error}
\begin{enumerate}[label=(\roman*)]
\setlength{\itemsep}{0pt}
\item Generate bootstrap samples $X_1^{\ast b},\ldots,X_n^{\ast b}$ for $b=1,\ldots,B$.
\item Compute $\widehat{\psi}^{\ast b}$ on each bootstrap sample.
\item Estimate standard error by the sample standard deviation of $\widehat{\psi}^{\ast1},\ldots,\widehat{\psi}^{\ast B}$.
\end{enumerate}
\end{method}

\begin{definition} \label{def:parametric-bootstrap} \hlt{Parametric Bootstrap}\\
In a parametric model, the \hlt{parametric bootstrap} first fits $\widehat{\theta}$ and then simulates pseudo-data from $P_{\widehat{\theta}}$.  It uses the fitted model rather than the empirical distribution as the resampling law.
\end{definition}

\begin{definition} \label{def:permutation-test} \hlt{Permutation Test}\\
A \hlt{permutation test} uses a group of transformations $G$ under which the null distribution is invariant.  For a test statistic $T$, the permutation reference distribution is
\[
\{T(gX):g\in G\},
\]
or a random sample from this set when $G$ is large.
\end{definition}

\begin{remark} \label{rem:resampling-validity} \hlt{Validity Depends on the Resampling Law}\\
Bootstrap validity is an asymptotic approximation statement; permutation validity is often finite-sample exact under invariance.  Both require the resampling mechanism to mimic the sampling variation relevant to the target statistic.
\end{remark}
